% SPDX-License-Identifier: CC-BY-4.0
% Copyright (c) 2026 Martin Schröder <info@swedishembedded.com>

\section{The coding-task experiment: protocol}\label{sec:protocol}

The question that matters to a team, once the loop of
\cref{sec:architecture,sec:admission,sec:verify,sec:dataset} exists, is:
\emph{given a git repository, can a local agentic model be tuned to work competently in
it, repeatedly and with further repositories, and does it get better at what it has met
rather than only at the newest?} This section states the protocol that asks it. It was
registered before any run (\code{prereg/coding-tasks.md}; its dated changes are listed in \cref{app:deviations}); everything
in it is a registered rule or a statement about the code; the coding-task measurements are
reported in \cref{sec:pending}.\cl{PR00}

\paragraph{Tasks.} A task is a snapshot of a repository tree at a revision, an
instruction stating a symptom, and privileged material the solver never receives: the
hidden checks (fail-to-pass and pass-to-pass commands) and the paths that carry them.\cl{PR01}
A candidate task is admitted only if the snapshot's selected tests pass twice, a change makes
a previously passing test fail reproducibly, and the reverse change restores green
(\cref{fig:valid}); a red or flaky baseline stops the run.\cl{PR01,PS01} Splinter's repository
tasks are made by code and never proposed by a model, and have their own refusal
vocabulary.\cl{A01,A09} The failing tests are moved out of the tree the solver receives, and a
repair passes when the hidden files are overlaid fresh on the patched tree and the
fail-to-pass and pass-to-pass commands pass; a patch that touches a protected path
fails.\cl{V07}

\begin{figure}[t]
\centering
\begin{tikzpicture}[
  stp/.style={draw,rounded corners,align=center,minimum width=2.2cm,minimum height=0.9cm,font=\scriptsize},
  arr/.style={-{Stealth[length=1.6mm]},thick},
  rej/.style={draw,dashed,rounded corners,align=center,font=\scriptsize,text=red!60!black}]
\node[stp] (a) {Baseline\\passes $\times2$};
\node[stp,right=0.55cm of a] (b) {Apply\\change};
\node[stp,right=0.55cm of b] (c) {Some test\\fails};
\node[stp,right=0.55cm of c] (d) {Failure\\reproducible};
\node[stp,right=0.55cm of d] (e) {Reverse patch\\restores green};
\node[stp,right=0.55cm of e,fill=green!12] (f) {Seal hidden\\files};
\foreach \x/\y in {a/b,b/c,c/d,d/e,e/f} \draw[arr] (\x) -- (\y);
\node[rej,below=0.5cm of a] (ra) {flaky or red:\\refuse};
\node[rej,below=0.5cm of c] (rc) {no failing test\\(equivalent): refuse};
\node[rej,below=0.5cm of e] (re) {not restorable:\\refuse};
\draw[arr,dashed,red!60!black] (a) -- (ra);
\draw[arr,dashed,red!60!black] (c) -- (rc);
\draw[arr,dashed,red!60!black] (e) -- (re);
\end{tikzpicture}
\caption{Admission of a mutation-seeded repository task. Every refusal is counted by reason.}\label{fig:valid}
\end{figure}

\paragraph{Solving, recording and trajectory data.} The solver is a Sven coding worker with
file, search and shell tools, a tool-call budget and a repetition stop, started in a fresh
copy of the snapshot. One harness profile is pinned for every solve of every arm and for data
collection alike: at most two follow-up turns, at most sixty tool
calls and a repetition stop at three; a follow-up is given when the visible checks stay red
and also when a turn ends with the tree unchanged, asking the worker to make the change by
calling the edit tools.\cl{R03,PS01,PR14} Measurements made before this profile are not
comparable and are excluded from every analysis; the reason is a pilot observation recorded in
the lane notes, in which the untrained 8B model almost never edited a file under the earlier
unbounded harness.\cl{PR14,PR15} The
recorded unit is the full trajectory; the verdict is produced outside the agent. Only
attempts whose patch passes and whose exposure audit is clean are eligible for training, and
for the headline experiment the attempts are those of a stronger teacher model, not of the
student (\cref{sec:pending});\cl{PR15}\cl{L06,L08} each assistant action becomes a training record conditioned on the prior
conversation and on the tool definitions the agent was offered, with system and user turns
carrying no loss, and compacted trajectories are excluded rather than repaired.\cl{D04,D05,D06,AU02}

\paragraph{Validity and leakage.} Four registered rules keep the exam honest. (i)~Families are
assigned to partitions by a stable hash before any rollout, and a family of the qualification
partition is never solved for data purposes.\cl{PR02,L07} (ii)~Hidden files are absent from the
solver's tree and are materialised only in the verifier's.\cl{V07} (iii)~An exposure audit fails
an attempt whose tool calls name a hidden path or whose outputs repeat a shingle of hidden
content, and exposed attempts and harness failures are withheld and counted in no
denominator of the paired test; in the release gate's improvement check they count as tasks
the arm did not solve, so a candidate cannot remove a loss from the test.\cl{L06,L08,PR06} (iv)~The isolation level of every run is recorded; with
\code{none} a host command is still confined by the kernel (\cref{sec:audit}) and the audit
is a second line, and a task whose soundness needs a container is not admitted.\cl{PR06,AU09,CN01} Whether a repository is in the base model's
pretraining data is not known; the base arm on the same exam bounds that effect without
removing it.

\paragraph{Units, arms and the paired exam.} The unit of analysis is the family, the file a
fault was seeded in for the authored repositories and the subsystem for real
ones.\cl{PR02} Tasks within a family are not independent, so no test treats them as such. The
base arm is the unmodified checkpoint in a pinned harness profile; the adapter arm is the same
checkpoint with the adapter attached; system prompt, tool schemas, budgets, decoding (the model family's recommended sampling, each task solved under three registered seeds and pass@1 averaged over them, thinking off as the primary variant and a 1024-token thinking budget as the secondary one) and environment image are identical.\cl{PR03,PR35} A family is won by the adapter
when its resolved fraction is higher, lost when lower, and dropped when tied. The primary
decision is a one-sided exact sign test over the non-tied qualification families at
$\alpha=\alphalevel$, on the adapter served through Brain and driven by Sven; a result that is
not significant is \emph{not shown}, never equivalence.\cl{PR04,G03}

\begin{table}[t]
\centering\small
\input{generated/design_table}
\caption{The exact one-sided sign test at $\alpha=\alphalevel$: wins required and power if
each non-tied family is a win with the stated probability (a design calculation computed by
\code{scripts/build_data.py}, not a measurement).}\label{tab:design}
\end{table}

\paragraph{What the test can decide.} \Cref{tab:design} and \cref{fig:power} give the wins
needed and the power at assumed win probabilities: with fewer than five non-tied families no
decision is possible, \mDesignNOneFiveWins{} wins are needed at fifteen non-tied families
and \mDesignNTwoFiveWins{} at twenty-five, where the power at a win share of 0.8 is
\mDesignNTwoFivePowerEightZero{}.\cl{CP01,CP02} The registration therefore asks for about
twenty-five non-tied qualification families, registered as 46 families (28 Python, 18 Lua) before any run, with the base-versus-teacher
discordance on development families reported as the sizing input, and labels a smaller comparison a pilot without changing the
rule.\cl{PR04}

\begin{figure}[t]
\centering
\begin{tikzpicture}
\begin{axis}[width=0.75\linewidth,height=5.2cm,xlabel={non-tied qualification families $n$},
  ylabel={power},ymin=0,ymax=1.02,xmin=2,xmax=40,legend pos=south east,
  legend style={font=\scriptsize},grid=major,
  every axis plot/.append style={thick}]
\addplot[blue,mark=none] table[col sep=comma,x=n,y=q70] {generated/power_curve.csv};
\addplot[red!70!black,mark=none,dashed] table[col sep=comma,x=n,y=q80] {generated/power_curve.csv};
\addplot[green!50!black,mark=none,dotted] table[col sep=comma,x=n,y=q90] {generated/power_curve.csv};
\legend{$P(\text{win})=0.7$,$P(\text{win})=0.8$,$P(\text{win})=0.9$}
\end{axis}
\end{tikzpicture}
\caption{Power of the exact one-sided sign test ($\alpha=\alphalevel$) against the number of
non-tied families. The sawtooth is the discreteness of the test.}\label{fig:power}
\end{figure}

\paragraph{Retention across rounds.} After each round $k$, the exam families of every earlier
round are solved again by the round-$k$ adapter and by its predecessor. Retention holds only if
the one-sided 95\% lower bound of the change in resolve rate (family-cluster bootstrap) exceeds
$-0.05$ absolute; a failed check refuses the candidate.\cl{PR05} The gate of
\cref{sec:gate} already holds each earlier release to the exam of its own run for repository
runs;\cl{G07} the gate's suite does not yet test acquired capabilities in every domain
(\cref{sec:requirements}), which is why the registered check is stricter
than the gate's.

\paragraph{Controls and stop rules.} A gain must survive three controls, each a separate
training run on the same budget: training on trajectories regardless of verdict (does
verification matter), an adapter trained on a different repository (did it learn the tool
protocol rather than the repository), and no replay of earlier rounds; for teacher-distilled data these are teacher trajectories
regardless of verdict and verified teacher trajectories of a different repository.\cl{PR07,PR15} Registered
stop rules include dropping a task kind whose verifier false-pass rate exceeds 5\% and reporting
a development gain without a qualification gain as memorisation.\cl{PR08}

\paragraph{The code that exists.} \code{splinter practice} implements the stages of a
single-repository round from a revision to a gated release.\cl{PS01} Eight authored repository
families ship with the repository (\cref{tab:families}); a family's tests must pass five times
identically and its faults be caught before it is catalogued.\cl{PS02} Two verbs apply the partitions and the paired comparison of this section to the base and the
served adapter.
\code{splinter tasks pick SET\ldots{} --family-share F --seed S} divides the union of task sets
into a picked set and the rest by a stable hash of the seed and the family, whole families at a
time (\code{--per-family N} caps each family instead, for the one-task-per-family verdict
sample). \code{splinter experiences compare CANDIDATE BASELINE} pairs the two arms' experiences
by task, counts a family once, and reports the exact one-sided sign test the release gate uses,
the per-family wins, losses and ties, quality, cost and autonomy with family-bootstrap
intervals, and the share of tasks on which the arms' verdicts agree; exposed and harness-failed
attempts are left out and counted.\cl{PS09} \Cref{sec:replicate} marks each step of the replication as available or planned.\cl{PL01}
